\documentclass[tikz,border=8pt]{standalone}
\usepackage{tikz}
\usepackage{amsmath}
\usepackage{amssymb}
\usepackage{pgfplots}
\pgfplotsset{compat=1.18}
\usepackage{ctex}
\usetikzlibrary{calc,positioning,arrows.meta,matrix}

% LC 295 Find Median from Data Stream
% 双堆（max-heap 左半 + min-heap 右半）3 快照演示
% 插入顺序：1, 2, 3, 4, 5

\begin{document}
\begin{tikzpicture}[
  every node/.style={font=\small},
  maxnode/.style={draw=orange!70, circle,
    minimum size=0.75cm, thick, fill=orange!10},
  minnode/.style={draw=blue!60, circle,
    minimum size=0.75cm, thick, fill=blue!8},
  newnode/.style={draw=green!60!black, circle,
    minimum size=0.75cm, very thick, fill=green!10},
  medbox/.style={draw=red!60, fill=red!8, rounded corners=3pt,
    font=\small\bfseries, inner sep=5pt},
  arr/.style={->,>=Stealth,gray!60,thick},
  labelbox/.style={draw=gray!40,fill=white,thin,rounded corners=2pt,
    font=\scriptsize,inner sep=4pt},
]

%% ── 标题 ──────────────────────────────────────────────────
\node[font=\large\bfseries] at (9.5, 1.2)
  {数据流中位数（LC 295 Find Median）};
\node[font=\small,gray] at (9.5, 0.55)
  {max-heap（左半小数）+ min-heap（右半大数）；两堆大小差 $\leq 1$};

%% ── 双堆结构说明 ─────────────────────────────────────────
\node[font=\scriptsize,orange!70,anchor=east] at (8.5, -0.3)
  {max-heap（左半）};
\node[font=\scriptsize,blue!70,anchor=west] at (10.5, -0.3)
  {min-heap（右半）};
\node[font=\scriptsize,gray] at (9.5, -0.3) {$|$};

%% ══════════════════════════════════════════════════════════
%% 快照 1：插入 1,2,3 后
%% max-heap: [3]（左半最大=3）
%% min-heap: [4,5]（右半最小=4，即待合并的索引号）
%% 实际上插入1,2,3：
%%   插1 → max:[1], min:[]  → 中位数 1
%%   插2 → max:[1], min:[2] → 中位数 (1+2)/2=1.5
%%   插3 → max:[2], min:[3] → 中位数 2（平衡后）
%%   插4 → max:[2], min:[3,4] → 中位数 (2+3)/2=2.5
%%   插5 → max:[3], min:[4,5] → 中位数 3
%% ══════════════════════════════════════════════════════════

\def\snapY{-1.6}

%% ── 快照 1：插入 1,2,3 后，中位数=2 ──────────────────────
\def\sAx{2.8}
\node[font=\scriptsize\bfseries,blue!70] at (\sAx+0.5, \snapY-0.1)
  {快照 1：插入 1, 2, 3};
\node[font=\tiny,gray] at (\sAx+0.5, \snapY-0.5) {中位数 = 2};

% max-heap：[2]（顶=2）
\node[maxnode] (ma0) at (\sAx-0.3, \snapY-1.5) {2};
\node[font=\tiny,orange!60,anchor=south] at (ma0.north) {max-top};

% min-heap：[3]（顶=3）
\node[minnode] (mi0) at (\sAx+1.3, \snapY-1.5) {3};
\node[font=\tiny,blue!60,anchor=south] at (mi0.north) {min-top};

% 中位数框
\node[medbox] at (\sAx+0.5, \snapY-2.8) {中位数 = 2};
\node[font=\tiny,gray] at (\sAx+0.5, \snapY-3.4) {max.top() = 2};

% 堆标签框
\node[labelbox,align=left] at (\sAx+0.5, \snapY-4.5)
  {max-heap: [2]\\min-heap: [3]\\大小: 1 vs 1};

%% ── 快照 2：插入 4 后，中位数=2.5 ──────────────────────
\def\sBx{8.0}
\node[font=\scriptsize\bfseries,blue!70] at (\sBx+1.0, \snapY-0.1)
  {快照 2：再插入 4};
\node[font=\tiny,gray] at (\sBx+1.0, \snapY-0.5) {中位数 = 2.5};

% max-heap：[2]（顶=2）
\node[maxnode] (mb0) at (\sBx+0.0, \snapY-1.5) {2};
\node[font=\tiny,orange!60,anchor=south] at (mb0.north) {max-top};

% min-heap：[3,4]（树形，3在顶）
\node[minnode] (mi1) at (\sBx+2.0, \snapY-1.5) {3};
\node[minnode] (mi2) at (\sBx+3.2, \snapY-2.6) {4};
\node[font=\tiny,blue!60,anchor=south] at (mi1.north) {min-top};

\draw[arr] (mi1) -- (mi2);

% 中位数框
\node[medbox] at (\sBx+1.0, \snapY-2.8) {中位数 = 2.5};
\node[font=\tiny,gray] at (\sBx+1.0, \snapY-3.4) {(max.top()+min.top())/2};

% 堆标签框
\node[labelbox,align=left] at (\sBx+1.0, \snapY-4.5)
  {max-heap: [2]\\min-heap: [3,4]\\大小: 1 vs 2};

%% ── 快照 3：插入 5 后，中位数=3，堆再平衡 ──────────────
\def\sCx{14.5}
\node[font=\scriptsize\bfseries,blue!70] at (\sCx+1.2, \snapY-0.1)
  {快照 3：再插入 5（平衡）};
\node[font=\tiny,gray] at (\sCx+1.2, \snapY-0.5) {中位数 = 3};

% max-heap：[3]（min-heap 顶 3 被移到 max-heap 平衡）
\node[maxnode] (mc0) at (\sCx+0.0, \snapY-1.5) {3};
\node[font=\tiny,orange!60,anchor=south] at (mc0.north) {max-top};

% min-heap：[4,5]（树形，4在顶）
\node[minnode] (mi3) at (\sCx+2.4, \snapY-1.5) {4};
\node[minnode] (mi4) at (\sCx+3.6, \snapY-2.6) {5};
\node[font=\tiny,blue!60,anchor=south] at (mi3.north) {min-top};
\draw[arr] (mi3) -- (mi4);

% 中位数框
\node[medbox] at (\sCx+1.2, \snapY-2.8) {中位数 = 3};
\node[font=\tiny,gray] at (\sCx+1.2, \snapY-3.4) {max.top() = 3（奇数）};

% 堆标签框
\node[labelbox,align=left] at (\sCx+1.2, \snapY-4.5)
  {max-heap: [3]\\min-heap: [4,5]\\大小: 1 vs 2（$\to$均衡: 2 vs 2后再push 5: 1 vs 2）};

% 注释：平衡操作
\node[labelbox,align=left,anchor=north west] at (\sCx-0.5, \snapY-5.5)
  {平衡：min.top(3) 移入 max-heap\\max-heap: [3], min-heap: [4,5]};

%% ── 数据流插入流程（完整序列）────────────────────────────
\node[font=\scriptsize,gray,anchor=west] at (0.5, -8.5)
  {插入序列：};
\foreach \i/\v/\med in {
  1/1/1.0,
  2/2/1.5,
  3/3/2.0,
  4/4/2.5,
  5/5/3.0
} {
  \node[draw=gray!40,fill=white,thin,rounded corners=2pt,
        font=\scriptsize,inner sep=4pt,align=center,minimum width=2.2cm]
    at (1.5+\i*2.8, -8.5)
    {插入 \v\\中位数=\med};
}

%% ── 分隔线 ──────────────────────────────────────────────
\draw[gray!25,dashed] (0.0, -9.7) -- (20.0, -9.7);

%% ── 算法说明 ─────────────────────────────────────────────
\node[draw=gray!50,fill=white,rounded corners=3pt,font=\small,
      inner sep=8pt,align=left]
  at (9.5, -10.8)
  {\textbf{双堆中位数算法（addNum: $O(\log n)$，findMedian: $O(1)$）：}\\[4pt]
   维护两个堆：\texttt{max\_heap}（存较小的一半）、\texttt{min\_heap}（存较大的一半）\\[2pt]
   \textbf{插入 x：}先 push 到 max\_heap，再将 max\_heap.top() 移到 min\_heap\\
   \quad\quad\quad\quad 若 min\_heap.size() > max\_heap.size()，再将 min\_heap.top() 移回 max\_heap\\[2pt]
   \textbf{求中位数：}两堆等大 $\to$ 两顶平均；max\_heap 更大 $\to$ max\_heap.top()};

\end{tikzpicture}
\end{document}
